\documentclass[a4paper,11pt]{article}
\usepackage[utf8]{inputenc}
\usepackage[T1]{fontenc}
\usepackage{amsmath,amssymb,amsthm}
\usepackage{algorithm}
\usepackage{algpseudocode}
\usepackage{listings}
\usepackage{xcolor}
\usepackage{hyperref}
\usepackage{url}
\usepackage{graphicx}
\usepackage{booktabs}
\usepackage{array}
\usepackage{longtable}
\usepackage{geometry}
\usepackage{natbib}

\geometry{margin=2.5cm}

%% Python listing style
\lstset{
  language=Python,
  basicstyle=\small\ttfamily,
  keywordstyle=\color{blue}\bfseries,
  commentstyle=\color{gray}\itshape,
  stringstyle=\color{red!60!black},
  showstringspaces=false,
  breaklines=true,
  frame=single,
  numbers=left,
  numberstyle=\tiny\color{gray},
  numbersep=5pt,
  xleftmargin=15pt,
}

\newtheorem{definition}{Definition}
\newtheorem{example}{Example}

\title{%
  \textbf{pateda: A Python Library for Estimation of Distribution Algorithms}\\[0.5em]
  \large User Guide --- Version 1.0 (Draft)
}
\author{
  Roberto Santana\\
  Department of Computer Science and Artificial Intelligence\\
  University of the Basque Country (UPV/EHU)\\
  \texttt{roberto.santana@ehu.eus}
}
\date{2026}

\begin{document}
\maketitle

\begin{abstract}
This document describes \textbf{pateda}, a Python library for estimation of
distribution algorithms (EDAs). pateda supports single- and multi-objective
optimization over discrete and continuous search spaces. The library provides
(i) a set of plug-and-play algorithm classes that encapsulate complete EDA
pipelines, (ii) a flexible component-based architecture that allows the user to
mix and match any combination of selection, learning, sampling, replacement,
and other operators, and (iii) a knowledge-extraction module for analyzing
probabilistic models learned during the evolutionary search.  pateda is the
Python successor to MATEDA-2.0 \citep{Santana_et_al:2009b}, inheriting its
modular philosophy while taking advantage of modern Python scientific computing
libraries.
\end{abstract}

\tableofcontents
\newpage

%% ============================================================
\section{Introduction}
%% ============================================================

Estimation of distribution algorithms (EDAs)
\citep{Larranhaga_and_Lozano:2002,Lozano_et_al:2005,Muhlenbein_and_Paas:1996,Pelikan_et_al:2006}
are a class of evolutionary algorithms that replace the variation operators of
traditional genetic algorithms (GAs) \citep{Goldberg:1989,Holland:1975} with
the explicit estimation and sampling of probabilistic models.  At each
generation, a probabilistic model is learned from a selected subset of the
current population; new individuals are then generated by sampling from this
model.  The probabilistic model captures the statistical dependencies between
the problem variables that tend to co-occur in high-quality solutions, endowing
EDAs with the ability to exploit problem structure in a principled way.

Several classes of probabilistic models have been proposed for EDAs.  They
range from simple univariate (independent-variable) models
\citep{Muhlenbein_and_Paas:1996,Baluja:1994} through tree-structured models
\citep{Baluja_and_Davies:1997,Pelikan_and_Muhlenbein:1999} and Bayesian
networks \citep{Etxeberria_and_Larranhaga:1999,Pelikan_et_al:1999,Pelikan:2005}
to Markov networks
\citep{Santana:2003c,Santana:2005,Shakya_and_McCall:2007,Shakya_and_Santana:2008},
factorized distributions \citep{Muhlenbein_et_al:1999,Harik:1999}, and mixtures
of distributions \citep{Bosman_and_Thierens:2002}.  For continuous optimization,
Gaussian EDAs
\citep{Larranhaga_et_al:1999,Larranhaga_et_al:2000,Bosman_and_Thierens:2000,Bosman_and_Thierens:2000a}
and mixture-of-Gaussian models are the most common choices.

Beyond optimization, the probabilistic models learned by EDAs encode information
about the problem structure that can be extracted and analyzed.  This
\emph{knowledge extraction} perspective has been pursued in a number of works
\citep{Echegoyen_et_al:2007,Echegoyen_et_al:2008,Hauschild_and_Pelikan:2008,Lima_et_al:2007,Santana_et_al:2007b}
and is an important feature of pateda.

\textbf{pateda} is a Python~3 library that implements a wide range of EDAs for
discrete and continuous single- and multi-objective optimization.  It is the
Python successor to MATEDA-2.0 \citep{Santana_et_al:2009b}, a Matlab suite
with a similar modular philosophy.  The main design goals of pateda are:

\begin{itemize}
  \item \textbf{Plug-and-play convenience.}  Algorithm classes such as
    \texttt{UMDA}, \texttt{EBNA}, \texttt{GaussianUMDA}, and
    \texttt{GaussianNetworkEDA} provide a one-liner interface for the most
    common EDA configurations.

  \item \textbf{Flexible component assembly.}  The \texttt{EDAComponents}
    dataclass allows the user to compose any EDA from independently swappable
    components: seeding, selection, learning, sampling, replacement, mutation,
    crossover, local optimization, repairing, statistics, and stopping condition.

  \item \textbf{Extensibility.}  Every component type is defined as an abstract
    base class.  Adding a new learning algorithm, selection scheme, or test
    function requires only the implementation of the corresponding abstract
    method.

  \item \textbf{Single and multi-objective support.}  The fitness function may
    return a scalar or a vector.  Multi-objective selection (Pareto
    non-dominated sorting, crowding distance, indicator-based), decomposition
    (MOEA/D), and archive management are provided.

  \item \textbf{Knowledge extraction.}  A dedicated module provides tools for
    analyzing the structure of the learned models, computing dependency
    measures, and visualizing network structures.
\end{itemize}

The remainder of this guide is organized as follows.
Section~\ref{sec:install} describes installation and the package structure.
Section~\ref{sec:eda} explains how to define and execute a generic EDA using
the core \texttt{EDA} class and \texttt{EDAComponents}.
Section~\ref{sec:algorithms} surveys the plug-and-play algorithm classes.
Section~\ref{sec:models} details the probabilistic model families, organized
by algorithm family.
Section~\ref{sec:components} documents all EDA component types and their
implementations.
Section~\ref{sec:functions} lists the built-in benchmark functions and test problems.
Section~\ref{sec:multiobjective} covers the multi-objective infrastructure.
Section~\ref{sec:knowledge} describes the knowledge-extraction module.
Section~\ref{sec:apriori} explains mechanisms for injecting a-priori structural
information.

%% ============================================================
\section{Installation and Package Structure}
\label{sec:install}
%% ============================================================

\subsection{Installation}

pateda is distributed as a Python package located in the \texttt{packages/pateda}
subdirectory of the repository.  Install it in editable mode using:

\begin{lstlisting}[language=bash]
pip install -e packages/pateda
\end{lstlisting}

\noindent Required dependencies: \texttt{numpy}, \texttt{scipy},
\texttt{scikit-learn}, \texttt{networkx}, \texttt{matplotlib}.
Optional: \texttt{seaborn} (visualization).

\subsection{Package Structure}

After installation, the top-level \texttt{pateda} namespace provides direct
access to the most commonly used classes:

\begin{lstlisting}
from pateda import (
    EDA, EDAComponents,          # core
    UMDA, TreeEDA, EBNA, BOA,   # discrete plug-and-play
    MNFDA, MNFDAG, MOA, FDA,    # Markov network / FDA
    MKEDA, MTED, AffEDA,        # chain / mixture / affinity
    CUMDA, CFDA, PBIL, BSC,     # constrained / incremental
    GaussianUMDA, GaussianFull, # continuous
    GaussianNetworkEDA,          # Gaussian network
    MixtureGaussianEDA,          # mixture of Gaussians
)
\end{lstlisting}

The internal module tree is:

\begin{center}
\begin{tabular}{lp{9cm}}
\toprule
\textbf{Module} & \textbf{Contents} \\
\midrule
\texttt{core/} & \texttt{EDA} executor, \texttt{EDAComponents}, \texttt{Statistics}, \texttt{Cache}, abstract component interfaces, model containers \\
\texttt{algorithms/} & Plug-and-play wrappers: \texttt{discrete.py}, \texttt{continuous.py}, \texttt{base.py} \\
\texttt{learning/} & All learning method implementations \\
\texttt{sampling/} & All sampling method implementations \\
\texttt{selection/} & Selection operators \\
\texttt{replacement/} & Replacement operators \\
\texttt{seeding/} & Population initialization methods \\
\texttt{mutation/} & Mutation operators \\
\texttt{crossover/} & Crossover operators \\
\texttt{repairing/} & Constraint repair methods \\
\texttt{local\_optimization/} & Local search methods \\
\texttt{inference/} & MAP and k-MAP inference \\
\texttt{multiobjective/} & Dominance, crowding, MOEA/D, indicators, archive \\
\texttt{stop\_conditions/} & Stopping criteria \\
\texttt{statistics/} & Population statistics and tracker \\
\texttt{functions/} & Benchmark functions (discrete and continuous) \\
\texttt{knowledge\_extraction/} & Dependency analysis, network measures, visualizations \\
\bottomrule
\end{tabular}
\end{center}

\subsection{Running the Examples}

The \texttt{packages/pateda/examples/} directory contains runnable scripts for
all major algorithm families.  For example:

\begin{lstlisting}[language=bash]
cd packages/pateda/examples
python umda_onemax.py
python ebna_deceptive.py
python gaussian_umda_sphere.py
python tree_eda_ising.py
python multiobjective_approaches_demo.py
\end{lstlisting}

%% ============================================================
\section{Defining and Executing an EDA}
\label{sec:eda}
%% ============================================================

\subsection{General Description}

The main component of pateda is the \texttt{EDA} class, defined in
\texttt{core/eda.py}.  It implements the generic EDA loop described in
Algorithm~\ref{alg:EDA}, analogous to \texttt{RunEDA.m} in MATEDA-2.0.

\begin{algorithm}
\caption{Estimation of Distribution Algorithm}
\label{alg:EDA}
\begin{algorithmic}[1]
\State Set $t \leftarrow 0$
\Repeat
  \If{$t = 0$}
    \State Generate initial population $D_0$ using the \emph{seeding method}
    \State Apply \emph{repairing method} to $D_0$ (if required)
    \State Evaluate $D_0$ using the \emph{fitness function}
    \State Apply \emph{local optimization} to $D_0$ (if required)
  \Else
    \State Sample $D_\text{new}$ from the model using the \emph{sampling method}
    \State Apply \emph{mutation} to $D_\text{new}$ (if required)
    \State Apply \emph{repairing method} to $D_\text{new}$ (if required)
    \State Evaluate $D_\text{new}$ using the \emph{fitness function}
    \State Apply \emph{local optimization} to $D_\text{new}$ (if required)
    \State Form $D_t$ from $D_{t-1}$ and $D_\text{new}$ using the \emph{replacement method}
  \EndIf
  \State Select $D_t^S \subseteq D_t$ using the \emph{selection method}
  \State Optionally weight individuals in $D_t^S$ (\emph{selection weighting})
  \State Learn a probabilistic model $M_t$ from $D_t^S$ using the \emph{learning method}
  \State $t \leftarrow t + 1$
\Until{the \emph{stopping condition} is satisfied}
\end{algorithmic}
\end{algorithm}

\subsection{The \texttt{EDA} Class}

An EDA is instantiated with:

\begin{lstlisting}
from pateda.core.eda import EDA

eda = EDA(
    pop_size    = 300,
    n_vars      = 20,
    fitness_func = my_func,
    cardinality  = np.full(20, 2),   # binary variables
    components   = my_components,
    random_seed  = 42,
)
stats, cache = eda.run(verbose=True)
\end{lstlisting}

\subsubsection{Input Parameters}

\begin{description}
  \item[\texttt{pop\_size}] Population size (integer).
  \item[\texttt{n\_vars}] Number of variables (integer).
  \item[\texttt{fitness\_func}] Callable that maps a 1-D NumPy array (an
    individual) to a scalar (single objective) or a 1-D array (multi-objective).
    pateda maximizes fitness; for minimization, negate the function.
  \item[\texttt{cardinality}] For discrete problems: a 1-D integer array of
    length \texttt{n\_vars} giving the number of values each variable can take.
    For continuous problems: a 2-D array of shape
    \texttt{(2, n\_vars)} where row~0 contains lower bounds and row~1 upper bounds.
  \item[\texttt{components}] An \texttt{EDAComponents} object specifying all
    the EDA building blocks (see Section~\ref{sec:components}).
  \item[\texttt{random\_seed}] Integer seed passed to NumPy's
    \texttt{default\_rng}.  Pass \texttt{None} for non-reproducible runs.
  \item[\texttt{selection\_weighting}] How the selected individuals are weighted
    when passed to the learning method.  Options: \texttt{"uniform"} (default,
    equivalent to classic EDA), \texttt{"proportional"} (fitness proportional),
    \texttt{"boltzmann"} (Boltzmann distribution with temperature
    \texttt{weighting\_beta}).
\end{description}

\subsubsection{Output: \texttt{Statistics} and \texttt{Cache}}

\texttt{eda.run()} returns a pair \texttt{(stats, cache)}.

The \texttt{Statistics} object provides:
\begin{itemize}
  \item \texttt{stats.best\_fitness}: list of best fitness per generation.
  \item \texttt{stats.mean\_fitness}: list of mean fitness per generation.
  \item \texttt{stats.std\_fitness}: list of standard deviation per generation.
  \item \texttt{stats.best\_individual}: the best solution found overall.
  \item \texttt{stats.best\_fitness\_overall}: its fitness value.
  \item \texttt{stats.generation\_found}: the generation at which it was found.
\end{itemize}

The \texttt{Cache} object stores per-generation data controlled by a
\texttt{CacheConfig}:
\begin{itemize}
  \item \texttt{cache.populations}: list of population arrays.
  \item \texttt{cache.fitness\_values}: list of fitness arrays.
  \item \texttt{cache.models}: list of learned models.
  \item \texttt{cache.selected\_populations}: list of selected subsets.
\end{itemize}

\subsection{The \texttt{EDAComponents} Dataclass}

The \texttt{EDAComponents} dataclass (defined in \texttt{core/components.py})
specifies all building blocks of the EDA.  All fields except
\texttt{learning}, \texttt{sampling}, and \texttt{stop\_condition} are optional.

\begin{lstlisting}
from pateda.core.components import EDAComponents
from pateda.seeding.random_init import RandomInit
from pateda.selection.truncation import TruncationSelection
from pateda.learning.umda import LearnUMDA
from pateda.sampling.fda import SampleFDA
from pateda.replacement.elitist import ElitistReplacement
from pateda.stop_conditions.max_generations import MaxGenerations

components = EDAComponents(
    seeding         = RandomInit(),
    selection       = TruncationSelection(ratio=0.5),
    learning        = LearnUMDA(alpha=0.0),
    sampling        = SampleFDA(n_samples=300),
    replacement     = ElitistReplacement(n_elite=1),
    stop_condition  = MaxGenerations(100),
)
\end{lstlisting}

\subsection{Example: UMDA on the OneMax Problem}
\label{ex:umda_onemax}

\begin{lstlisting}
import numpy as np
from pateda.core.eda import EDA
from pateda.core.components import EDAComponents
from pateda.seeding.random_init import RandomInit
from pateda.selection.truncation import TruncationSelection
from pateda.learning.umda import LearnUMDA
from pateda.sampling.fda import SampleFDA
from pateda.replacement.elitist import ElitistReplacement
from pateda.stop_conditions.max_generations_or_optimum import (
    MaxGenerationsOrOptimum)

n_vars   = 50
pop_size = 300

def onemax(x):
    return np.sum(x)

cardinality = np.full(n_vars, 2)   # binary variables

components = EDAComponents(
    seeding        = RandomInit(),
    selection      = TruncationSelection(ratio=0.5),
    learning       = LearnUMDA(alpha=0.0),
    sampling       = SampleFDA(n_samples=pop_size),
    replacement    = ElitistReplacement(n_elite=1),
    stop_condition = MaxGenerationsOrOptimum(
                         max_gen=200, optimum=float(n_vars)),
)

eda = EDA(pop_size, n_vars, onemax, cardinality,
          components, random_seed=42)
stats, cache = eda.run()
print("Best fitness:", stats.best_fitness_overall)
print("Best solution:", stats.best_individual)
\end{lstlisting}

\subsection{Example: Gaussian UMDA on the Sphere Function}
\label{ex:gumda_sphere}

\begin{lstlisting}
import numpy as np
from pateda import GaussianUMDA

def sphere(x):
    return -np.sum(x**2)    # maximise the negated sphere

alg = GaussianUMDA(
    n_vars      = 30,
    bounds      = (-5.0, 5.0),
    fitness_func= sphere,
    pop_size    = 500,
    n_gen       = 100,
    random_seed = 0,
)
stats, _ = alg.run()
print(f"Best: {stats.best_fitness_overall:.6f}")
\end{lstlisting}

\subsection{Fitness Convention}

pateda \textbf{maximizes} fitness.  When the objective is to be minimized (e.g.,
energy, error, distance), the user should negate the function:

\begin{lstlisting}
def minimization_objective(x):
    return -cost_function(x)   # pateda maximises this
\end{lstlisting}

%% ============================================================
\section{Plug-and-Play Algorithm Classes}
\label{sec:algorithms}
%% ============================================================

For convenience, pateda provides ready-to-use wrappers that assemble the
appropriate components internally.  These classes expose a minimal constructor
and a \texttt{run()} method that returns \texttt{(stats, cache)}.

\subsection{Discrete Algorithm Classes}

All discrete algorithm classes share the following common parameters:

\begin{center}
\begin{tabular}{lp{9cm}}
\toprule
Parameter & Description \\
\midrule
\texttt{n\_vars} & Number of variables \\
\texttt{cardinality} & \texttt{int} (same for all vars) or 1-D integer array \\
\texttt{fitness\_func} & Callable; higher is better \\
\texttt{pop\_size} & Population size (default: 100) \\
\texttt{n\_gen} & Number of generations (default: 50) \\
\texttt{selection\_ratio} & Truncation fraction (default: 0.5) \\
\texttt{elitism} & Preserve best individual (default: True) \\
\texttt{random\_seed} & RNG seed (default: None) \\
\texttt{alpha} & Laplace smoothing for probability tables (default: 0.0) \\
\bottomrule
\end{tabular}
\end{center}

\vspace{0.5em}

Table~\ref{tab:discrete_algs} summarizes the discrete algorithm classes.

\begin{table}[ht]
\centering
\small
\setlength{\tabcolsep}{4pt}
\begin{tabular}{llp{6.5cm}}
\toprule
\textbf{Class} & \textbf{Identifier} & \textbf{Description} \\
\midrule
\texttt{UMDA} & UMDA & Univariate marginal distribution algorithm
  \citep{Muhlenbein_and_Paas:1996} \\
\texttt{PBIL} & PBIL & Population-based incremental learning \citep{Baluja:1994} \\
\texttt{BSC}  & BSC  & Fitness-weighted univariate marginals (Bisection scheme \citep{Inza_et_al:2000}) \\
\texttt{MIMIC} & MIMIC & Mutual-information maximization for input clustering \citep{DeBonet_et_al:1996} \\
\texttt{BMDA} & BMDA & Bivariate marginal distribution algorithm \\
\texttt{TreeEDA} & Tree-EDA & Maximum-weight spanning tree of pairwise MI \citep{Baluja_and_Davies:1997,Pelikan_and_Muhlenbein:1999} \\
\texttt{TreeEDAR} & Tree-EDA\_r & Tree-EDA with restricted interaction matrix \\
\texttt{AffEDA} & AffEDA & Affinity-propagation cliques \citep{Frey_and_Dueck:2007,Santana_et_al:2008} \\
\texttt{MKEDA} & MK-EDA$k$ & $k$-order Markov chain EDA \citep{Santana_et_al:2004} \\
\texttt{MTED} & MT-EDA & Mixture of trees EDA \citep{Santana_et_al:2001br} \\
\texttt{MNFDA} & MN-FDA & Markov network FDA sampling \citep{Santana:2005} \\
\texttt{MNFDAR} & MN-FDA\_r & MN-FDA with random variable ordering \\
\texttt{MNFDAG} & MN-FDAg & Augmented Markov network + Gibbs sampling \\
\texttt{MNFDAGR} & MN-FDAg\_r & MN-FDAg with random Gibbs ordering \\
\texttt{MOA}  & MOA & Markovianity-based optimization algorithm \citep{Shakya_and_McCall:2007,Shakya_and_Santana:2008} \\
\texttt{EBNA} & EBNA & Estimation of Bayesian network algorithm \citep{Etxeberria_and_Larranhaga:1999} \\
\texttt{BOA}  & BOA & Bayesian optimization algorithm \citep{Pelikan_et_al:1999,Pelikan:2005} \\
\texttt{FDA}  & FDA & Factorized distribution algorithm \citep{Muhlenbein_et_al:1999} \\
\texttt{CUMDA} & CUMDA & Unitation-constrained UMDA \\
\texttt{CFDA} & CFDA & Unitation-constrained FDA \\
\bottomrule
\end{tabular}
\caption{Discrete plug-and-play algorithm classes in pateda.}
\label{tab:discrete_algs}
\end{table}

\subsection{Continuous Algorithm Classes}

Table~\ref{tab:continuous_algs} summarizes the continuous algorithm classes.

\begin{table}[ht]
\centering
\small
\setlength{\tabcolsep}{4pt}
\begin{tabular}{llp{7.5cm}}
\toprule
\textbf{Class} & \textbf{Identifier} & \textbf{Description} \\
\midrule
\texttt{GaussianUMDA} & G-UMDA & Univariate Gaussian marginals \citep{Larranhaga_et_al:1999,Larranhaga_et_al:2000} \\
\texttt{GaussianFull}  & G-Full & Full multivariate Gaussian \citep{Larranhaga_et_al:2000,Bosman_and_Thierens:2000} \\
\texttt{GaussianNetworkEDA} & G-Net & Gaussian network (GMRF) EDA \\
\texttt{MixtureGaussianEDA} & MGaussian & Mixture of Gaussians \citep{Bosman_and_Thierens:2002} \\
\bottomrule
\end{tabular}
\caption{Continuous plug-and-play algorithm classes in pateda.}
\label{tab:continuous_algs}
\end{table}

Common parameters for continuous classes are \texttt{n\_vars},
\texttt{bounds} (a tuple \texttt{(lo, hi)} or a \texttt{(2, n\_vars)} array),
\texttt{fitness\_func}, \texttt{pop\_size}, \texttt{n\_gen},
\texttt{selection\_ratio}, \texttt{elitism}, and \texttt{random\_seed}.

\subsection{Quick Comparison Example}

\begin{lstlisting}
from pateda import UMDA, TreeEDA, EBNA
import numpy as np

def deceptive3(x):
    """Sum of order-3 deceptive functions."""
    n = len(x)
    score = 0.0
    for i in range(0, n - 2, 3):
        u = int(x[i]) + int(x[i+1]) + int(x[i+2])
        score += 0.9 if u < 3 else 1.0
    return score

n_vars = 30
for AlgClass in [UMDA, TreeEDA, EBNA]:
    alg = AlgClass(n_vars=n_vars, cardinality=2,
                   fitness_func=deceptive3,
                   pop_size=500, n_gen=100, random_seed=1)
    stats, _ = alg.run(verbose=False)
    print(f"{AlgClass.__name__:10s}  best={stats.best_fitness_overall:.3f}")
\end{lstlisting}

%% ============================================================
\section{Probabilistic Models and Algorithm Families}
\label{sec:models}
%% ============================================================

\subsection{Discrete Representation}

Discrete variables take values in $\{0, 1, \ldots, r_i - 1\}$ where $r_i$ is
the cardinality of variable $X_i$.  Binary variables have $r_i = 2$.  pateda
represents a population of $N$ solutions as a 2-D NumPy array of shape
$(N, n)$ with integer dtype.  The \texttt{cardinality} parameter is a 1-D
integer array $[r_1, r_2, \ldots, r_n]$.  Mixed-cardinality problems (some
variables binary, others multi-valued) are fully supported.

\subsection{Continuous Representation}

Continuous variables are real-valued, bounded by
$[\ell_i, u_i]$.  The population is a 2-D float array of shape $(N, n)$.
The \texttt{cardinality} array for continuous problems is a 2-D array of
shape $(2, n)$: row~0 contains the lower bounds and row~1 the upper bounds.

\subsection{Factorized Distributions}
\label{sec:factorized}

A factorized distribution \citep{Muhlenbein_et_al:1999,Harik:1999} represents
a joint probability over $\mathbf{X} = (X_1, \ldots, X_n)$ as a product of
factor (clique) marginals:

\begin{equation}
  p(\mathbf{x}) = \prod_{i=1}^{m} p(x_{C_i} \mid x_{S_i})
\end{equation}

where $C_i$ is the set of ``new'' variables in factor $i$ and $S_i$ is the
set of ``separator'' variables shared with previous factors.  This class of
models, called \emph{ordered factorizations} \citep{Santana:2005}, subsumes:
\begin{itemize}
  \item Marginal product factorizations (univariate, bivariate, etc.)
  \item Markov chain factorizations
  \item Junction-tree factorizations
  \item Junction-graph factorizations \citep{Santana_et_al:2008a}
\end{itemize}

pateda stores factorizations in the \texttt{FactorizedModel} container, which
holds a list of factors; each factor records its clique variables, separator
variables, and a conditional probability table.  The \texttt{LearnFDA} class
fits the tables from a given population, and \texttt{SampleFDA} samples new
solutions using probabilistic logic sampling (PLS) \citep{Muhlenbein_and_Paas:1996}.

\textbf{FDA} (\texttt{FDA} class): when a user supplies a fixed clique
structure (known from the problem domain), \texttt{LearnFDA} computes only
the conditional probability tables from data, bypassing structure learning
entirely.  This is the factorized distribution algorithm of
\citet{Muhlenbein_et_al:1999}.

\subsection{Univariate Models}
\label{sec:univariate}

Univariate models assume all variables are independent: $p(\mathbf{x}) =
\prod_i p(x_i)$.  The marginal $p(x_i)$ is estimated as the empirical
frequency in the selected population (with optional Laplace smoothing
$\alpha$).

\begin{description}
  \item[\textbf{UMDA}] \citep{Muhlenbein_and_Paas:1996}.
    Univariate marginal distribution algorithm.  Each marginal is learned
    independently from the selected population.
    Learning: \texttt{LearnUMDA}. Sampling: \texttt{SampleFDA}.

  \item[\textbf{PBIL}] \citep{Baluja:1994}.
    Population-based incremental learning.  Maintains a running estimate
    of the marginals updated with learning rate $\alpha$:
    $P^{(t+1)}(x_i) = (1 - \alpha)\,P^{(t)}(x_i) + \alpha\,\hat{P}(x_i)$.
    With $\alpha = 1$ PBIL reduces to UMDA.
    Learning: \texttt{LearnPBIL}. Sampling: \texttt{SampleFDA}.

  \item[\textbf{BSC}] \citep{Inza_et_al:2000}.
    Bisection-based fitness-weighted marginals.
    $P(x_i = k) \propto \sum_{\{j: x_i^{(j)} = k\}} f(x^{(j)})$.
    Learning: \texttt{LearnBSC}. Sampling: \texttt{SampleFDA}.

  \item[\textbf{MIMIC}] \citep{DeBonet_et_al:1996}.
    Mutual Information Maximization for Input Clustering.
    Learns a chain-structured factorization that maximizes the sum of
    pairwise mutual information along the chain.
    Learning: \texttt{LearnMIMIC}. Sampling: \texttt{SampleFDA}.
\end{description}

\subsection{Tree and Forest Models}
\label{sec:trees}

Tree-structured models represent pairwise variable dependencies using a
maximum-weight spanning tree of the mutual-information matrix \citep{Chow_and_Liu:1968}.

\begin{description}
  \item[\textbf{TreeEDA / Tree-EDA}] \citep{Baluja_and_Davies:1997,Pelikan_and_Muhlenbein:1999}.
    Learns a directed tree using Chow--Liu algorithm on pairwise MI.
    Learning: \texttt{LearnTreeModel}. Sampling: \texttt{SampleFDA}.

  \item[\textbf{TreeEDAR / Tree-EDA\_r}]
    Variant of Tree-EDA with a user-supplied \texttt{interaction\_matrix} that
    constrains which pairs of variables can form a dependency.
    Learning: \texttt{LearnTreeModelR}.

  \item[\textbf{BMDA}]
    Bivariate marginal distribution algorithm.  Detects pairwise dependencies
    with chi-square tests and builds a forest of bivariate conditionals.
    Learning: \texttt{LearnBMDA}.

  \item[\textbf{AffEDA}]
    Affinity-propagation EDA \citep{Frey_and_Dueck:2007,Santana_et_al:2008}.
    Uses affinity propagation on the MI matrix to discover variable cliques,
    then builds a factorized model over those cliques.
    Learning: \texttt{LearnAffinityFactorization}.

  \item[\textbf{MTED / MT-EDA}] \citep{Santana_et_al:2001br}.
    Mixture of trees EDA.  Learns $k$ tree components with mixture weights.
    Weight learning options: \texttt{"uniform"}, \texttt{"em"}, or
    \texttt{"fitness\_proportional"}.  Adaptive and prior-based extensions
    are supported to prevent premature convergence.
    Learning: \texttt{LearnMixtureTrees}. Sampling: \texttt{SampleMixtureTrees}.
\end{description}

\subsection{Bayesian Network Models}
\label{sec:bayes}

Bayesian networks \citep{Lauritzen_and_Spiegelhalter:1988} are directed acyclic
graphs in which $X_i$ is conditionally independent of its non-descendants given
its parents $\mathbf{Pa}_i$.  The joint distribution factorizes as:

\begin{equation}
  p(\mathbf{x}) = \prod_{i=1}^{n} p(x_i \mid \mathbf{pa}_i).
\end{equation}

The local distribution $p(x_i \mid \mathbf{pa}_i)$ is a conditional probability
table (CPT) with $r_i$ entries for each of the $q_i = \prod_{X_j \in
\mathbf{Pa}_i} r_j$ parent configurations.

\begin{description}
  \item[\textbf{EBNA}] \citep{Etxeberria_and_Larranhaga:1999}.
    Estimation of Bayesian network algorithm.  Structure learning uses a
    greedy score-and-search procedure with the Bayesian information criterion
    (BIC) \citep{Schwarz:1978} or the Bayesian score
    \citep{Cooper_and_Herskovits:1992}.
    Learning: \texttt{LearnEBNA}. Sampling: \texttt{SampleBayesianNetwork}.

  \item[\textbf{BOA}] \citep{Pelikan_et_al:1999,Pelikan:2005}.
    Bayesian optimization algorithm.  Greedy structure search with BIC/MDL
    scoring.
    Learning: \texttt{LearnBOA}. Sampling: \texttt{SampleBayesianNetwork}.
\end{description}

Both EBNA and BOA sample new individuals using probabilistic logic sampling
over the directed network.  Additionally, pateda provides sampling strategies
that find the \emph{most probable explanation} (MPE): the configuration
$\mathbf{x}^* = \arg\max_\mathbf{x} p(\mathbf{x})$, computed via
\texttt{SampleGibbs} or the inference methods in \texttt{inference/map\_inference.py}.

\subsection{Markov Network Models}
\label{sec:markov}

A Markov random field (MRF) \citep{Hammersley_and_Clifford:1968} is an
undirected graphical model in which the joint distribution factorizes over the
cliques of the graph as a Gibbs distribution:

\begin{equation}
  p(\mathbf{x}) = \frac{1}{Z}\exp\!\left(-H(\mathbf{x})\right),\quad
  H(\mathbf{x}) = \sum_{C \in \mathcal{C}} \Phi_C(\mathbf{x}_C)
\end{equation}

where $\mathcal{C}$ is the set of maximal cliques and $Z$ is the partition
function.  Sampling from MRFs typically uses Gibbs sampling: iteratively
sample each variable $X_i$ from its conditional distribution given the
current values of its neighbors $N(X_i)$.

\begin{description}
  \item[\textbf{MN-FDA / MNFDA}] \citep{Santana:2005}.
    Learns a Markov network structure from chi-square independence tests
    between pairs of variables and represents the joint distribution as
    a factorized model.  Sampling uses PLS (factorized distribution sampling).
    Learning: \texttt{LearnMNFDA}. Sampling: \texttt{SampleFDA}.

  \item[\textbf{MN-FDA\_r / MNFDAR}]
    Variant of MN-FDA with random variable ordering during sampling, improving
    exploration of multimodal distributions.
    Learning: \texttt{LearnMNFDAR}. Sampling: \texttt{SampleFDA}.

  \item[\textbf{MN-FDAg / MNFDAG}]
    Augmented Markov network EDA.  Extends MN-FDA with additional edges
    based on mutual information.  Sampling uses Gibbs sampling.
    Learning: \texttt{LearnMNFDAG}. Sampling: \texttt{SampleGibbs}.

  \item[\textbf{MN-FDAg\_r / MNFDAGR}]
    MN-FDAg with random Gibbs sampling order.
    Learning: \texttt{LearnMNFDAGR}. Sampling: \texttt{SampleGibbs}.

  \item[\textbf{MOA}] \citep{Shakya_and_McCall:2007,Shakya_and_Santana:2008}.
    Markovianity-based optimization algorithm.  Learns local $k$-neighborhoods
    for each variable and samples via Gibbs sampling.
    Learning: \texttt{LearnMOA}. Sampling: \texttt{SampleGibbs}.
\end{description}

\subsection{Markov Chain Models}
\label{sec:markov_chain}

\begin{description}
  \item[\textbf{MK-EDA / MKEDA}] \citep{Santana_et_al:2004}.
    $k$-order Markov chain EDA.  Each variable $X_i$ depends on the $k$
    preceding variables in a fixed linear ordering:
    $p(\mathbf{x}) = p(x_1)\prod_{i=2}^{n} p(x_i \mid x_{i-k}, \ldots, x_{i-1})$.
    Well-suited to problems with a natural sequential structure (e.g.\
    protein folding with relative encodings).
    Learning: \texttt{LearnMarkovChain}. Sampling: \texttt{SampleMarkovChain}.
\end{description}

\subsection{Constrained EDAs}
\label{sec:constrained}

\begin{description}
  \item[\textbf{CUMDA}]
    Unitation-constrained UMDA.  Every solution must contain exactly
    \texttt{n\_ones} ones.  The seeding method
    (\texttt{SeedingUnitationConstraint}) generates feasible initial
    solutions, and the sampling method (\texttt{SampleCUMDA}) uses
    Stochastic Universal Sampling to maintain the constraint.
    Learning: \texttt{LearnCUMDA}.

  \item[\textbf{CFDA}]
    Unitation-constrained FDA.  Uses a tree-structured factorized model
    while enforcing the same fixed-unitation constraint as CUMDA.
    Learning: \texttt{LearnCFDA}. Sampling: \texttt{SampleCFDA}.
\end{description}

\subsection{Continuous Gaussian Models}
\label{sec:gaussian}

\begin{description}
  \item[\textbf{GaussianUMDA}] \citep{Larranhaga_et_al:1999,Larranhaga_et_al:2000}.
    Each variable is modeled independently as a Gaussian:
    $p(x_i) = \mathcal{N}(x_i; \mu_i, \sigma_i^2)$.
    Parameters $(\mu_i, \sigma_i^2)$ are estimated from the selected population.
    Learning: \texttt{LearnGaussianUnivariate}. Sampling: \texttt{SampleGaussianUnivariate}.

  \item[\textbf{GaussianFull}] \citep{Larranhaga_et_al:2000,Bosman_and_Thierens:2000}.
    Full multivariate Gaussian: $p(\mathbf{x}) = \mathcal{N}(\mathbf{x};\boldsymbol{\mu},\Sigma)$.
    All pairwise dependencies are captured by the covariance matrix $\Sigma$.
    May suffer from variance collapse in later generations
    \citep{Bosman_and_Grahl:2008}; a variance scaling parameter is provided.
    Learning: \texttt{LearnGaussianFull}. Sampling: \texttt{SampleGaussianFull}.

  \item[\textbf{GaussianNetworkEDA}]
    Gaussian network (GMRF) EDA.  Learns a sparse Gaussian graphical model
    where each variable's local distribution is a linear regression on its
    parents in a directed graph.  The GMRF approach identifies relevant
    conditional independences in continuous data.

  \item[\textbf{MixtureGaussianEDA}] \citep{Bosman_and_Thierens:2002}.
    Mixture of Gaussian distributions.  The population is first clustered
    (k-means or similar), then a separate Gaussian (univariate or full) is
    fit to each cluster.  Suitable for multimodal search landscapes.
    Learning: \texttt{LearnMixtureGaussian}. Sampling: \texttt{SampleMixtureGaussian}.
\end{description}

%% ============================================================
\section{EDA Components}
\label{sec:components}
%% ============================================================

\subsection{Seeding Methods}

Seeding methods initialize the population at generation 0.

\begin{description}
  \item[\texttt{RandomInit}]
    Samples each variable uniformly from its range.
    For discrete variables: integer in $[0, r_i)$.
    For continuous variables: real in $[\ell_i, u_i]$.

  \item[\texttt{BiasInit}]
    Biased binary initialization.  Variable $X_i$ is set to 1 with probability
    $p$ and to 0 with probability $1-p$.

  \item[\texttt{SeedingUnitationConstraint}]
    Generates binary solutions with exactly \texttt{num\_ones} ones.
    Required by \texttt{CUMDA} and \texttt{CFDA}.

  \item[\texttt{SeedThisPop}]
    Seeds the population with a user-supplied array.
\end{description}

\subsection{Selection Methods}

Selection determines which individuals are used for model learning.
The role of selection in EDAs has been studied by
\citet{Johnson_and_Shapiro:2002,Lima_et_al:2007,Mahnig_and_Muehlenbein:2001,Santana_et_al:2005b}.

\begin{description}
  \item[\texttt{TruncationSelection}]
    Selects the top $\lfloor \texttt{ratio} \cdot N \rfloor$ individuals by
    fitness.  This is the most common selection method in EDAs.

  \item[\texttt{TournamentSelection}]
    Tournament selection of configurable size.

  \item[\texttt{ProportionalSelection} / \texttt{SUSSelection}]
    Fitness-proportional selection using stochastic universal sampling (SUS)
    \citep{Baker:1987}.

  \item[\texttt{BoltzmannSelection}]
    Selection with Boltzmann temperature; probability proportional to
    $e^{\beta f(x)}$.

  \item[\texttt{RankingSelection}]
    Rank-based selection; probability proportional to $2r/(N(N+1))$ where
    $r$ is the fitness rank.

  \item[\texttt{NonDominatedSelection}]
    Selects the non-dominated front.  For multi-objective problems.

  \item[\texttt{CrowdingSelection}]
    NSGA-II-style: non-dominated sorting + crowding distance tie-breaking.

  \item[\texttt{IndicatorBasedSelection}]
    SMS-EMOA-style hypervolume indicator-based selection.

  \item[\texttt{ParetoFrontSelection}]
    Selects complete Pareto fronts up to the selection quota.
\end{description}

\subsection{Replacement Methods}

\begin{description}
  \item[\texttt{ElitistReplacement}]
    Copies the best \texttt{n\_elite} solutions from the previous generation
    into the new population if they are better than the worst new individual.

  \item[\texttt{GenerationalReplacement}]
    New population completely replaces the old one.
\end{description}

\subsection{Stopping Conditions}

\begin{description}
  \item[\texttt{MaxGenerations}]
    Stops after a fixed number of generations.

  \item[\texttt{MaxGenerationsOrOptimum}]
    Stops when either the generation limit is reached or a target fitness
    value is attained.
\end{description}

\subsection{Local Optimization}

\begin{description}
  \item[\texttt{GreedySearch}]
    Greedy hill climbing in the continuous case.

  \item[\texttt{DiscreteGreedySearch}]
    Greedy flip-based local search for discrete problems.

  \item[\texttt{DiscreteSimulatedAnnealing}]
    Simulated annealing for discrete problems.

  \item[\texttt{ContiguousBlockOptimization}]
    Specialized local search for contiguous-block problems.

  \item[\texttt{ScipyLocalSearch}]
    Wraps \texttt{scipy.optimize.minimize} for continuous local optimization.
\end{description}

\subsection{Repairing Methods}

\begin{description}
  \item[\texttt{BoundsRepair}]
    Clips continuous variable values to their feasible bounds.

  \item[\texttt{TrigonometricRepair}]
    Trigonometric projection for angular variables.

  \item[\texttt{UnitationRepair}]
    Adjusts binary solutions to satisfy a fixed unitation constraint.
\end{description}

%% ============================================================
\section{Test Functions and Benchmark Problems}
\label{sec:functions}
%% ============================================================

\subsection{Discrete Functions}

\begin{description}
  \item[\textbf{OneMax}] Maximizes $f(\mathbf{x}) = \sum_i x_i$.
    The simplest EDA benchmark; all variables are independent.

  \item[\textbf{Trap functions}] Deceptive order-$k$ subfunctions
    \citep{Ackley:1987}:
    $f(\mathbf{x}) = \sum_{j} \text{trap}_k(u_j)$ where $u_j$ is the
    number of ones in block $j$.  Optimal solution is the all-ones string;
    the deceptive attractor is the all-zeros string.

  \item[\textbf{Deceptive functions}] Goldberg's deceptive functions
    \citep{Goldberg:1987}.

  \item[\textbf{Additively decomposable functions (ADFs)}]
    $f(\mathbf{x}) = \sum_i f_i(\mathbf{x}_{S_i})$ where each subfunction
    $f_i$ is defined over a disjoint subset $S_i$ of variables.  Ideal for
    studying the structural identification capabilities of EDAs.

  \item[\textbf{NK landscapes}] \citep{Kauffman:1993}.
    Each variable depends on $k$ others drawn uniformly or in circular topology.
    The subfunction value for each configuration is drawn at random.

  \item[\textbf{Ising model}]
    Energy minimization on a lattice:
    $E = -\sum_{i<j} J_{ij}\sigma_i\sigma_j - \sum_i h_i\sigma_i$.

  \item[\textbf{HP protein folding}] \citep{Dill:1985}.
    2-D/3-D lattice model.  Solutions encode the sequence of relative moves of
    amino acid residues.  pateda uses the relative encoding with a backtracking
    repair heuristic.

  \item[\textbf{SAT}] \citep{Selman_et_al:1992}.
    Satisfiability problems; single- and multi-objective versions available.

  \item[\textbf{UBQP}]
    Unconstrained binary quadratic programming:
    $f(\mathbf{x}) = \mathbf{x}^T Q \mathbf{x}$.

  \item[\textbf{Contiguous-block functions}]
    Reward solutions with long contiguous blocks of identical bits.

  \item[\textbf{Integer functions}]
    Functions over multi-valued (non-binary) discrete variables.
\end{description}

\subsection{Continuous Functions}

\begin{description}
  \item[\textbf{Standard benchmarks}] (\texttt{functions/continuous/benchmarks.py}):
    Sphere, Rastrigin, Rosenbrock, Ackley, Griewank.

  \item[\textbf{AB off-lattice protein model}] \citep{Larranhaga_et_al:1999}.
    Off-lattice protein model with Lennard-Jones interactions.
    Variables are bond angles in $[0, 2\pi]$.

  \item[\textbf{GNBG instances}]
    Generalized numerical benchmark generator instances
    (\texttt{GNBG\_Instances.Python-main/}).
\end{description}

%% ============================================================
\section{Multi-Objective Optimization}
\label{sec:multiobjective}
%% ============================================================

\subsection{Multi-Objective Problem Definition}

A multi-objective optimization problem (MOP) \citep{Deb:2001} requires
simultaneously optimizing $m$ objectives $f_1(\mathbf{x}), \ldots, f_m(\mathbf{x})$.
Since objectives usually conflict, the goal is to find the \emph{Pareto front}:
the set of non-dominated solutions where no single objective can be improved
without degrading another.

pateda supports multi-objective problems by returning a vector from the fitness
function:

\begin{lstlisting}
def bi_objective(x):
    f1 = np.sum(x)            # first objective (maximise)
    f2 = -np.sum((x - 0.5)**2)  # second objective (maximise)
    return np.array([f1, f2])
\end{lstlisting}

\subsection{Pareto Dominance}

pateda implements Pareto dominance in \texttt{multiobjective/dominance.py}.
Solution $\mathbf{x}$ dominates $\mathbf{y}$ if it is at least as good on
all objectives and strictly better on at least one.

\subsection{Multi-Objective Selection}

\begin{description}
  \item[\textbf{Non-dominated sorting + crowding (NSGA-II)}]
    \citep{Deb_et_al:2002a}.
    Solutions are ranked into Pareto fronts.  Within each front, crowding
    distance breaks ties.  Available via \texttt{CrowdingSelection}.

  \item[\textbf{Indicator-based (SMS-EMOA)}]
    Hypervolume contribution per individual.  Available via
    \texttt{IndicatorBasedSelection} and \texttt{multiobjective/indicators.py}.

  \item[\textbf{Decomposition (MOEA/D)}] \citep{Muhlenbein_et_al:1999}.
    Decomposes the MOP into scalar subproblems via weight vectors.
    pateda's MOEA/D is implemented in \texttt{multiobjective/moead.py}.
\end{description}

\subsection{Multi-Objective EDA Pipeline}

Below is a sketch of a multi-objective EDA using NSGA-II-style selection with
a Bayesian network learner:

\begin{lstlisting}
from pateda.core.eda import EDA
from pateda.core.components import EDAComponents
from pateda.seeding.random_init import RandomInit
from pateda.selection.crowding import CrowdingSelection
from pateda.learning.ebna import LearnEBNA
from pateda.sampling.bayesian_network import SampleBayesianNetwork
from pateda.replacement.elitist import ElitistReplacement
from pateda.stop_conditions.max_generations import MaxGenerations
import numpy as np

n_vars, pop_size = 20, 500

def mo_func(x):
    f1 = np.sum(x)
    f2 = np.sum(1 - x)
    return np.array([f1, f2])

components = EDAComponents(
    seeding        = RandomInit(),
    selection      = CrowdingSelection(n_select=pop_size // 2),
    learning       = LearnEBNA(),
    sampling       = SampleBayesianNetwork(n_samples=pop_size),
    replacement    = ElitistReplacement(n_elite=1),
    stop_condition = MaxGenerations(50),
)
eda = EDA(pop_size, n_vars, mo_func,
          np.full(n_vars, 2), components)
stats, cache = eda.run()
\end{lstlisting}

\subsection{Extracting Pareto Fronts}

After running the EDA with \texttt{cache\_config} that stores populations and
fitness values, the full Pareto set approximation can be extracted:

\begin{lstlisting}
from pateda.selection.utils.pareto import find_pareto_front
import numpy as np

all_pop = np.vstack(cache.populations)
all_fit = np.vstack(cache.fitness_values)
pareto_idx = find_pareto_front(all_fit)
pareto_pop = all_pop[pareto_idx]
pareto_fit = all_fit[pareto_idx]
\end{lstlisting}

%% ============================================================
\section{Knowledge Extraction}
\label{sec:knowledge}
%% ============================================================

One of the most distinctive features of pateda---inherited from MATEDA-2.0---is
the ability to extract and visualize knowledge from the probabilistic models
learned during the EDA search.  The \texttt{knowledge\_extraction} module
provides tools organized into the following categories.

\subsection{Dependency Analysis}

\texttt{knowledge\_extraction/dependency\_analysis.py} provides:

\begin{description}
  \item[\texttt{compute\_correlation\_matrix}]
    Computes the Pearson correlation matrix between variables from a population.

  \item[\texttt{learn\_bayesian\_network}]
    Learns a Bayesian network structure from a population using score-and-search.

  \item[\texttt{learn\_gaussian\_network}]
    Learns a Gaussian graphical model from continuous data.

  \item[\texttt{analyze\_variable\_dependencies}]
    Computes pairwise mutual information and tests for marginal independence.

  \item[\texttt{compute\_mutual\_information}]
    Mutual information between two discrete random variables.
\end{description}

\subsection{Fitness Measures}

\texttt{knowledge\_extraction/fitness\_measures.py} provides tools to
quantify the EDA's behavior over generations:

\begin{description}
  \item[Response to selection $R(t)$]
    $R(t) = \bar{f}(t+1) - \bar{f}(t)$: improvement in mean fitness
    from one generation to the next.

  \item[Amount of selection $S(t)$]
    $S(t) = \bar{f}_s(t) - \bar{f}(t)$: difference between the mean
    fitness of the selected subset and the whole population.

  \item[Realized heritability $b(t)$]
    $b(t) = R(t) / S(t)$: fraction of the selection differential that
    is realized in the next generation \citep{Muhlenbein:1997}.
\end{description}

\subsection{Network Measures and Visualizations}

For models based on graphs (Bayesian networks, Markov networks, trees),
pateda tracks which edges appear across generations and runs:

\begin{description}
  \item[\textbf{Edge frequency matrix}]
    \citep{Echegoyen_et_al:2007,Santana_et_al:2007b}.
    The $(i,j)$ entry counts how often the edge $(X_i, X_j)$ appeared
    in the structures learned across all runs and generations.  Bright
    colors in the heat-map indicate frequent structural relationships,
    often corresponding to true problem-structure interactions.

  \item[\textbf{Dendrogram of edges}]
    Hierarchical clustering of edges by co-occurrence frequency, displayed
    as a dendrogram \citep{Hauschild_and_Pelikan:2008}.

  \item[\textbf{Parallel coordinate view}]
    Each axis is a selected edge; each line connects edges that co-occur
    in the same structure at the same generation.

  \item[\textbf{Glyph representation}]
    A grid of glyphs (one per run $\times$ generation) showing which
    edges are present in each learned structure.
\end{description}

\subsection{Continuous Variable Analysis}

For continuous EDAs, \texttt{knowledge\_extraction/continuous\_visualizations.py}
and \texttt{gaussian\_networks.py} provide tools to plot the evolution of
Gaussian parameters (means, variances, correlations) across generations.
The \texttt{vine\_analysis.py} module supports analysis of vine copula structures.

%% ============================================================
\section{Injecting A-Priori Knowledge}
\label{sec:apriori}
%% ============================================================

A key advantage of EDA-based approaches is that a-priori information about
the problem structure can be directly incorporated into the model, potentially
reducing the search space and accelerating convergence
\citep{Santana_et_al:2009b}.

\subsection{Fixed Factorizations}

When the structure of the optimal model is known a priori (e.g., from domain
knowledge or problem decomposition), the user can supply a fixed clique
structure to \texttt{FDA}:

\begin{lstlisting}
import numpy as np
from pateda import FDA

# Fixed Markov-chain structure for a 10-variable problem
n = 10
cliques = np.array([
    [0, 1, 0],       # factor 1: only X_1
    [1, 1, 0, 1],    # factor 2: X_1 | X_2
    [1, 1, 1, 2],    # factor 3: X_2 | X_3
    # ... etc.
])

alg = FDA(n_vars=n, cardinality=2,
          fitness_func=my_func,
          cliques=cliques,
          pop_size=300, n_gen=50)
stats, _ = alg.run()
\end{lstlisting}

The Markov-chain FDA with order 1 is the natural model for problems with
sequential structure, such as the HP protein folding model
\citep{Dill:1985,Santana_et_al:2004}.

\subsection{Interaction Matrices}

\texttt{TreeEDAR}, \texttt{MNFDAR}, and \texttt{MNFDAGR} accept an
\texttt{interaction\_matrix} parameter: an $n \times n$ binary matrix where
entry $(i,j)=1$ means that variables $X_i$ and $X_j$ may be connected in
the learned model, and $(i,j)=0$ forbids the connection.  This allows the
user to encode known independence relationships or restrict the model to
biologically or physically plausible interactions.

\subsection{Biased Initialization}

\texttt{BiasInit} seeds the initial population with variable-wise biases.
For example, if domain knowledge indicates that variable $X_5$ is likely to
take value 1 in good solutions, a higher prior probability can be assigned:

\begin{lstlisting}
from pateda.seeding.bias_init import BiasInit

seeder = BiasInit(probs=np.array([0.5]*20 + [0.9]))
# Variable 21 is initialized with P(X=1) = 0.9
\end{lstlisting}

\subsection{Constrained Search}

\texttt{CUMDA} and \texttt{CFDA} enforce a fixed number of ones in every
solution.  This is useful for combinatorial problems where a cardinality
constraint must be satisfied throughout the search (e.g., feature selection
with exactly $k$ features).

\subsection{Selection Weighting}

The \texttt{selection\_weighting} parameter of the \texttt{EDA} constructor
allows fitness-proportional or Boltzmann weighting of the selected individuals
when computing model parameters, effectively blending structure learning with
fitness guidance:

\begin{lstlisting}
eda = EDA(pop_size, n_vars, fitness_func, cardinality,
          components,
          selection_weighting="boltzmann",
          weighting_beta=2.0)   # higher beta → stronger fitness pressure
\end{lstlisting}

%% ============================================================
\section{Conclusions}
%% ============================================================

pateda is a comprehensive Python library for estimation of distribution
algorithms supporting:

\begin{itemize}
  \item Discrete and continuous optimization, single and multi-objective.
  \item A broad range of EDA families: univariate, tree, Bayesian network,
    Markov network, factorized, mixture, Gaussian, constrained.
  \item Flexible component assembly via \texttt{EDAComponents}, enabling
    rapid prototyping of novel EDA variants.
  \item Mechanisms for incorporating a-priori problem structure knowledge.
  \item A knowledge-extraction module for model analysis and visualization.
  \item A rich library of benchmark functions including classical combinatorial
    problems, continuous benchmarks, and real-world problem instances.
\end{itemize}

Future versions will extend the guide with detailed worked examples for each
algorithm family, additional multi-objective benchmarks (DTLZ, ZDT),
performance comparison tables, and deeper documentation of the knowledge-
extraction and visualization tools.

%% ============================================================
\bibliographystyle{apalike}
\bibliography{pateda}
%% ============================================================

\end{document}
